\begin{center}
	\large{\textbf{\textcolor{black}{Integraci�n de subsistemas para el prototipo de EAMMRA.}}}\\
	\vspace{1pc}
	\textcolor {black}{Mariano CINQUINI, Patricio BOS, Igor PRARIO y Rui MARQUES ROJO}
\end{center}

\bigskip
	\begin{center}
		\textcolor{black}{\textbf{RESUMEN}}	
	\end{center}
	\begin*
		\indent
		\textit{En este informe t�cnico se presenta la integraci�n de cinco subsistemas para el prototipo de  Estaci�n Aut�noma Mar�tima para el Monitoreo de Ruido Ambiente (EAMMRA), en el marco del proyecto de I+D PIDDEF 32/14 del Ministerio de Defensa.  Los subsistemas que se integran son: alimentaci�n, control, procesamiento, adquisici�n y sensores. Se documentan sus principales caracter�sticas y se describen las interfaces f�sicas y l�gicas entre subsistemas. Asimismo, se hace una descripci�n funcional del ciclo de trabajo de la estaci�n a cargo de los subsistemas de control y procesamiento. Finalmente, se incluyen pruebas preliminares de laboratorio para evaluar el desempe�o del sistema integrado en su conjunto. }
	\end*



\bigskip
	\begin{center}
		\textcolor{black}{\textbf{ABSTRACT}}
		\end{center}
			\begin*
			\indent
		\textit{In this technical report the integration of five subsystems for the prototype of the Autonomous Maritime Station for Environmental Noise Monitoring (EAMMRA) is presented, within the framework of the R\&D project PIDDEF 32/14 of the Ministry of Defense. The integrated subsystems are: power, control, processing, acquisition and sensors. Their main characteristics are documented and the physical and logical interfaces between subsystems are described. Likewise, a functional description of the work cycle of the station in charge of the control and processing subsystems is made. Finally, preliminary laboratory tests are included to evaluate the performance of the integrated system as a whole.}        
	\end*

\clearpage